\section*{अध्याय \ref{ch:b2:angiosperm-reproduction} --- सपुष्पक पौधों का लैंगिक जनन}

\begin{solution}{exo:b2:angiosperm-reproduction:1}
बाह्यदल (बाह्यदलपुंज), दलपत्र (दलपुंज), \omterm{def:b2:angiosperm-reproduction:flower}{पुंकेसर} (पुमंग), \omterm{def:b2:angiosperm-reproduction:flower}{अंडप}
(जायांग)। \omterm{def:b2:angiosperm-reproduction:flower}{पुंकेसर}: तंतु और परागकोश। \omterm{def:b2:angiosperm-reproduction:flower}{अंडप}: \omterm{def:b2:angiosperm-reproduction:flower}{वर्तिकाग्र}, वर्तिका,
और \omterm{def:b2:plant-life-cycles:heterospory}{बीजांडों} सहित \omterm{def:b2:angiosperm-reproduction:flower}{अंडाशय}। ये सब \omterm{def:b2:meiosis-heredity:meiosis}{द्विगुणित} \omterm{def:b2:plant-life-cycles:cycles}{बीजाणुद्भिद} हैं; केवल \omterm{prop:b2:angiosperm-reproduction:pollen}{पराग
कण} (परागकोश के भीतर) और भ्रूणकोष (\omterm{def:b2:plant-life-cycles:heterospory}{बीजांड} के भीतर)
\omterm{def:b2:meiosis-heredity:meiosis}{अगुणित} \omterm{def:b2:plant-life-cycles:cycles}{युग्मकोद्भिद} हैं।
\end{solution}

\begin{solution}{exo:b2:angiosperm-reproduction:2}
\omterm{def:b2:plant-life-cycles:heterospory}{लघुबीजाणु} $n$; कायिक कोशिका $n$; शुक्राणु $n$; अंडाणु $n$; \omterm{prop:b2:angiosperm-reproduction:embryosac}{ध्रुवीय
केंद्रक} $n$; युग्मनज $2n$; भ्रूणपोष $3n$; \omterm{def:b2:angiosperm-reproduction:seed}{बीज-आवरण} $2n$ (मातृ
अध्यावरण); फल-भित्ति $2n$ (मातृ \omterm{def:b2:angiosperm-reproduction:flower}{अंडाशय})।
\end{solution}

\begin{solution}{exo:b2:angiosperm-reproduction:3}
वायु (घास, बलूत, भोजपत्र, हेज़ल, बरतंग): छोटे हरे \omterm{def:b2:angiosperm-reproduction:flower}{पुष्प}, न
दलपत्र, न गंध, न \omterm{def:b2:angiosperm-reproduction:pollination}{मकरंद}, लंबे तंतुओं पर झूलते परागकोश,
पंखनुमा \omterm{def:b2:angiosperm-reproduction:flower}{वर्तिकाग्र}, चिकने सूखे \omterm{def:b2:plant-life-cycles:heterospory}{पराग} की विपुल मात्रा, और पत्तियों से
पहले फूलना। मधुमक्खी (तिपतिया, सेज, फ़ॉक्सग्लव, लैवेंडर, सेब): रंगीन
दलपत्र, प्रायः पराबैंगनी संकेतों सहित, उतरने का चबूतरा, गंध,
\omterm{def:b2:angiosperm-reproduction:pollination}{मकरंद}, मध्यम मात्रा में चिपचिपा तराशा हुआ \omterm{def:b2:plant-life-cycles:heterospory}{पराग}, और दिन के उजाले में
फूलना।
\end{solution}

\begin{solution}{exo:b2:angiosperm-reproduction:4}
एक ही \omterm{thm:b2:angiosperm-reproduction:tube}{पराग नलिका} की दोनों शुक्राणु कोशिकाएँ संलयित होती हैं, एक
अंडाणु से (युग्मनज, $2n$, यानी भ्रूण) और एक केंद्रीय कोशिका से (भ्रूणपोष
केंद्रक, $3n$, यानी भोजन-भंडार)। भ्रूणपोष निषेचन के बाद ही बढ़ना आरंभ
करता है, इसलिए पौधा केवल उन्हीं \omterm{def:b2:plant-life-cycles:heterospory}{बीजांडों} को रसद देता है जिनमें कोई
भ्रूण है, जबकि अनावृतबीजी अपना भोजन-भंडार पहले ही गढ़ लेता है।
\end{solution}

\begin{solution}{exo:b2:angiosperm-reproduction:5}
$25/1.2 = \qty{21}{h}$। दिन 5, 6 और 7 को निकलते रेशम \omterm{def:b2:plant-life-cycles:heterospory}{पराग} पाते हैं;
दिन 8 से 12 वाले नहीं: रेशमों का $3/8$ परागित होता है, और
भुट्टी पाँच-आठवाँ भाग बंजर रहती है।
\end{solution}

\begin{solution}{exo:b2:angiosperm-reproduction:6}
$S_1S_2\times S_1S_2$: 0। $S_1S_2\times S_1S_3$: $\tfrac12$ (केवल $S_3$
\omterm{def:b2:plant-life-cycles:heterospory}{पराग} बढ़ता है)। $S_1S_2\times S_3S_4$: पूरा। दूसरे संकरण की संतति:
अंडाणु $S_1$ या $S_2$, शुक्राणु $S_3$: $S_1S_3$ और $S_2S_3$, आधे-आधे।
\end{solution}

\begin{solution}{exo:b2:angiosperm-reproduction:7}
कोई पौधा $n$ युग्मविकल्पियों में से 2 ढोता है और इनमें से कोई भी ढोता
\omterm{def:b2:plant-life-cycles:heterospory}{पराग} अस्वीकार कर देता है: योग्य अंश $(n - 2)/n$। किसी नए \omterm{def:b2:meiosis-heredity:vocabulary}{युग्मविकल्पी} को
कोई स्त्रीकेसर अस्वीकार नहीं करता, इसलिए उसका \omterm{def:b2:plant-life-cycles:heterospory}{पराग} औसत से अधिक \omterm{def:b2:angiosperm-reproduction:seed}{बीज} जनता
है और वह तब तक फैलता है जब तक बाक़ी जितना आम न हो जाए; और यही लाभ हर
विरल \omterm{def:b2:meiosis-heredity:vocabulary}{युग्मविकल्पी} को लोप से बचाता है। इसलिए वरण \omterm{def:b2:meiosis-heredity:vocabulary}{युग्मविकल्पी} जमा करता
जाता है, और स्व-अयोग्य जातियों की प्राकृतिक आबादियाँ दर्जनों ढोती हैं।
\end{solution}

\begin{solution}{exo:b2:angiosperm-reproduction:8}
दोनों \omterm{prop:b2:angiosperm-reproduction:embryosac}{ध्रुवीय केंद्रक} एक ही \omterm{def:b2:plant-life-cycles:heterospory}{गुरुबीजाणु} की प्रतियाँ हैं, इसलिए केंद्रीय
कोशिका $YY$ या $yy$ है (आधी-आधी); और शुक्राणु $Y$ या $y$ है (आधा-आधा):
भ्रूणपोष $YYY$, $YYy$, $Yyy$, $yyy$, हर एक $\tfrac14$; यानी तीन
चौथाई पीले। $yy$ से परागित होने पर: $YYy$ और $yyy$, आधे-आधे --- यानी
आधे दाने पीले।
\end{solution}

\begin{solution}{exo:b2:angiosperm-reproduction:9}
जौ के दाने बीच से काटे जाते हैं; भ्रूण वाले आधे भाग अपना मंड पचा लेते
हैं, और भ्रूण-रहित आधे भाग नहीं, जब तक उनमें जिबरेलिन न डाला जाए, और तब
वे पचा लेते हैं। (क) भ्रूण संकेत का स्रोत है, एंजाइम का नहीं।
(ख) जिबरेलिन ही संकेत है, और अपने आप में पर्याप्त।
(ग) ऐमिलेज़ भ्रूणपोष की ऐल्यूरोन परत बनाती है,
जो उस हॉर्मोन को उत्तर देती है। भ्रूण-रहित आधा भाग संकेत को
अनुक्रिया से अलग कर देता है: उसके बिना यह बताया ही न जाता कि मंड भ्रूण
स्वयं पचाता है या नहीं।
\end{solution}

\begin{solution}{exo:b2:angiosperm-reproduction:10}
वायु: प्रति \omterm{def:b2:plant-life-cycles:heterospory}{बीजांड} $10^{6}\times\qty{1}{\micro g} = \qty{1}{g}$। प्राणी:
प्रति \omterm{def:b2:plant-life-cycles:heterospory}{बीजांड} \qty{1}{mg} \omterm{def:b2:plant-life-cycles:heterospory}{पराग} जमा \qty{0.5}{mg} शर्करा $= \qty{1.5}{mg}$
--- यानी लगभग सात सौ गुना सस्ता। वायु-परागण
को कोई साथी नहीं चाहिए: वह कीटों के उड़ने से पहले आरंभिक वसंत में चलता
है, ठंडे, हवादार और खुले स्थानों में, रात में और वर्षा में, एक ही जाति
के सघन झुंडों में, और उसे न मकरंद-चोर ठग सकते हैं न कोई
विलुप्त हो जाने वाला परागक बीच में छोड़ सकता है।
\end{solution}

\begin{solution}{exo:b2:angiosperm-reproduction:11}
जो \omterm{def:b2:angiosperm-reproduction:seed}{बीज} जनक के पैरों में गिरते हैं वे उसी की छाया में, उसी की जड़ों के
बीच, और वहाँ उतरते हैं जहाँ उसी के विशिष्ट बीज-परभक्षी तथा रोगजनक जमा
रहते हैं; लगभग सारे मर जाते हैं। कोई पक्षी किसी \omterm{def:b2:angiosperm-reproduction:seed}{बीज} को कई किलोमीटर
(बीस मिनट में \qty{12}{km}) ले जाता है और उसे खाद सहित किसी
बाड़ या खुली जगह में गिरा देता है; और यदि सौ में एक \omterm{def:b2:angiosperm-reproduction:seed}{बीज} ही किसी अच्छी
जगह उतरे, तो भी वह वृक्ष के नीचे के सारे \omterm{def:b2:angiosperm-reproduction:seed}{बीजों} से अधिक है, और वह
आबादी गिरते \omterm{def:b2:angiosperm-reproduction:seed}{फल} की चाल से नहीं, पक्षियों की चाल से फैलती है। गूदा
टिकट का दाम है।
\end{solution}

\begin{solution}{exo:b2:angiosperm-reproduction:12}
घिरा होना \omterm{def:b2:plant-life-cycles:heterospory}{पराग} को मातृ ऊतक से होकर उगने पर मजबूर कर देता है, जो
स्त्रीकेसर को उसे परखने देता है (स्व-अयोग्यता, पराई जातियों की
अस्वीकृति), अनेक नलिकाओं को होड़ करने देता है ताकि सबसे तेज़ नलिका ही
\omterm{def:b2:angiosperm-reproduction:seed}{बीज} जने, \omterm{def:b2:plant-life-cycles:heterospory}{बीजांडों} को सूखने और शाकाहारियों से बचाता है, और
\omterm{def:b2:angiosperm-reproduction:flower}{अंडाशय} को ऐसा \omterm{def:b2:angiosperm-reproduction:seed}{फल} बना देता है जो प्रकीर्णन के लिए प्राणियों को भर्ती कर
ले; और फिर \omterm{prop:b2:angiosperm-reproduction:double}{द्विनिषेचन} केवल निषेचित \omterm{def:b2:plant-life-cycles:heterospory}{बीजांडों} को ही रसद देता है। क़ीमतें:
किसी \omterm{prop:b2:angiosperm-reproduction:pollen}{पराग कण} को \omterm{def:b2:plant-life-cycles:heterospory}{बीजांड} तक नहीं बल्कि किसी \omterm{def:b2:angiosperm-reproduction:flower}{वर्तिकाग्र} तक पहुँचना पड़ता
है, इसलिए एक वर्तिका, एक नलिका और उसका पथ-प्रदर्शन चाहिए, और \omterm{def:b2:angiosperm-reproduction:seed}{फल} एक भारी
निवेश है। यह संतुलन आवृतबीजियों के इतना पक्ष में रहा कि वे दस करोड़
वर्षों में शून्य से बढ़कर पादप जातियों के नौ दसवें भाग तक पहुँच गए।
\end{solution}

\begin{solution}{pb:b2:angiosperm-reproduction:1}
\textbf{1.} A ($S_1S_2$) पर: \omterm{def:b2:plant-life-cycles:heterospory}{पराग} A $f = 0$; B ($S_1$, $S_3$) $f =
\tfrac12$; C ($S_3$, $S_4$) $f = 1$।
\textbf{2.} B ($S_1S_3$) पर: A $\tfrac12$, B 0, C $\tfrac12$। C
($S_3S_4$) पर: A 1, B $\tfrac12$, C 0।
\textbf{3.} 50 योग्य कण; $0.1\times 50 = 5$ \omterm{def:b2:plant-life-cycles:heterospory}{बीजांड}।
\textbf{4.} C से: 100 योग्य, $0.1\times 100 = 10$, यानी सारे
दस \omterm{def:b2:plant-life-cycles:heterospory}{बीजांड}; A से: कोई नहीं।
\textbf{5.} C की यात्राओं से बने \omterm{def:b2:angiosperm-reproduction:seed}{फल} (10 \omterm{def:b2:angiosperm-reproduction:seed}{बीज}) रखे जाते हैं; B से (5 \omterm{def:b2:angiosperm-reproduction:seed}{बीज})
बस-बस रखे जाते हैं; A से गिरा दिए जाते हैं। केवल A का बग़ीचा कुछ नहीं देता।
\textbf{6.} आधे \omterm{def:b2:angiosperm-reproduction:flower}{पुष्प} \omterm{def:b2:angiosperm-reproduction:seed}{फल} बाँधते हैं (वे जिन्हें C से यात्रा मिली): प्रति
वृक्ष 2500 \omterm{def:b2:angiosperm-reproduction:seed}{फल}, यानी आवश्यक 250 का दस गुना --- पूरी फ़सल, और
वृक्ष अतिरिक्त गिरा देगा।
\textbf{7.} वह हर क़िस्म के पूरे पुष्पन-काल में योग्य \omterm{def:b2:plant-life-cycles:heterospory}{पराग} झाड़ता रहता है;
पर यदि वह किसी क़िस्म से दोनों $S$ \omterm{def:b2:meiosis-heredity:vocabulary}{युग्मविकल्पी} साझा करे
तो उस पर कुछ भी परागित नहीं करता।
\textbf{8.} प्रति वर्ग मीटर प्रति सेकंड $8\times 10^{7}/(7\times 86400) = 132$
कण।
\textbf{9.} प्रति घन मीटर $132/0.25 = 530$ कण।
\textbf{10.} प्रति सेकंड $132\times 10^{-4} = 0.013$: यानी हर
\qty{76}{s} में एक कण।
\textbf{11.} एक दिन में $0.013\times 86400 \approx 1100$ कण; और जो नलिका
पहले \omterm{def:b2:plant-life-cycles:heterospory}{बीजांड} तक पहुँचती है वही उसे निषेचित करती है, और फिर वह \omterm{def:b2:plant-life-cycles:heterospory}{बीजांड} बाक़ी
सबके लिए बंद हो जाता है।
\textbf{12.} \omterm{def:b2:angiosperm-reproduction:pollination}{परागण} के \qty{20}{h} बाद।
\textbf{13.} \omterm{def:b2:plant-life-cycles:heterospory}{पराग} दिन 0--7, रेशम दिन 5 से: केवल दिन 5, 6 और
7 अतिव्यापित होते हैं --- यानी रेशमों का $3/8$, और जो दिन 7 के बाद
निकलते हैं उन्हें कुछ नहीं मिलता; भुट्टी अधिकतर बंजर रहती है।
रेशम निकलने के समय सूखा मक्का की फ़सल के बिगड़ने का उत्कृष्ट कारण है।
\textbf{14.} भ्रूण $Yy$; भ्रूणपोष $YYy$।
\textbf{15.} पीला, क्योंकि $Y$ \omterm{def:b2:meiosis-heredity:vocabulary}{प्रभावी} है: पराग-जनक का लक्षण
माता-पौधे पर लगे \omterm{def:b2:angiosperm-reproduction:seed}{बीज} में दिखता है --- यही ज़ेनिया है।
\textbf{16.} केंद्रीय कोशिका $YY$ या $yy$: भ्रूणपोष $YYy$ (आधे) और
$yyy$ (आधे)।
\textbf{17.} $YYy$: 20 इकाई; $Yyy$: 10 इकाई; और ऐसा $yyY$ जिसका $Y$ पिता
से आया हो वही \omterm{def:b2:meiosis-heredity:vocabulary}{जीन-प्ररूप} $Yyy$ है, यानी 10 इकाई। दो पितृ
मात्राएँ हो ही नहीं सकतीं: केंद्रीय कोशिका को एक ही शुक्राणु निषेचित करता
है, और भ्रूणपोष में सदा दो मातृ जीनोम बनाम एक पितृ होते हैं।
\textbf{18.} वह मंड-भंडार है; दाने का $0.3/0.33 = \qty{91}{\%}$।
\textbf{19.} वह भंडार निषेचन के बाद ही गढ़ा जाता है, इसलिए उन \omterm{def:b2:plant-life-cycles:heterospory}{बीजांडों}
में कोई संसाधन नहीं जाता जिनमें कोई भ्रूण नहीं; जबकि अनावृतबीजी का
\omterm{def:b2:plant-life-cycles:cycles}{युग्मकोद्भिद} पहले ही गढ़ लिया जाता है और \omterm{def:b2:angiosperm-reproduction:pollination}{परागण} विफल हो तो व्यर्थ जाता है।
\textbf{20.} प्रति भुट्टी और प्रति पौधा $0.9\times 600 = 540$ दाने;
प्रति वर्ग मीटर $4320$; प्रति वर्ग मीटर $4320\times 0.3 = \qty{1.3}{kg}$
भ्रूणपोष।
\textbf{21.} \qty{13}{t/ha}।
\textbf{22.} दाने का कार्बन $0.5\times 1.2 = \qty{0.6}{kg/m^2}$;
भ्रूणपोष का कार्बन $0.45\times 1.3 = \qty{0.58}{kg/m^2}$: संगत
(छोटा शेष भ्रूण और आवरण है)।
\textbf{23.} प्रति दाना $8\times 10^{7}/4320 \approx \num{18500}$
कण।
\textbf{24.} उसका अपना पराग-बादल पतला है और उड़ जाता है, और उसकी
मंजरी अधिकतर उसके अपने रेशम निकलने से पहले ही झाड़ देती है, इसलिए अनेक
रेशम कभी परागित ही नहीं होते और भुट्टी में ख़ाली जगहें रह जाती हैं।
\textbf{25.} A खंड: आधे \omterm{def:b2:angiosperm-reproduction:flower}{पुष्प} \omterm{def:b2:angiosperm-reproduction:seed}{फल} बाँधते हैं, प्रति वृक्ष 2500, यानी
पूरी फ़सल। खेत: प्रति वर्ग मीटर 4320 दाने, \qty{13}{t/ha}।
\end{solution}
